% ============================================================
%  Dispensa 02f: Curva EIOPA Risk-Free — Metodo Smith-Wilson
%               Approccio variazionale (struttura stile PCA)
%  Corso di Calcolo Numerico — Laboratorio
%  Università degli Studi di Verona — A.A. 2026–2027
%
%  NOTA di approfondimento, non variante di lezione: isola e
%  sviluppa con molti più passaggi un solo argomento già
%  presente in forma compressa in 02c Sez. 4.2 (righe 789-919)
%  e in 02e Appendice A — il paragrafo "Semipositività: la
%  forma quadratica è un quadrato più l'integrale di un
%  quadrato". Nessuna sezione economica, nessuna calibrazione:
%  solo definizioni essenziali (richiamate, non dimostrate) +
%  la dimostrazione per esteso, con verifiche numeriche.
%  Nessuno script R né cartella output propri (nota solo
%  testuale/matematica).
% ============================================================
\documentclass[11pt,a4paper]{article}

\usepackage[T1]{fontenc}
\usepackage[utf8]{inputenc}
\usepackage[italian]{babel}
\usepackage{lmodern}
\usepackage{amsmath,amssymb,amsthm,mathtools}
\usepackage{geometry}
\usepackage{booktabs}
\usepackage{array}
\usepackage{enumitem}
\usepackage{graphicx}
\usepackage[table]{xcolor}
\usepackage{hyperref}
\usepackage{float}
\usepackage{fancyvrb}
\usepackage{tcolorbox}
\tcbuselibrary{breakable}
\usepackage{titlesec}

\graphicspath{{}}

\geometry{top=2.5cm,bottom=2.5cm,left=2.8cm,right=2.8cm}
\setlist[itemize]{topsep=4pt,itemsep=2pt,parsep=0pt}
\setlist[enumerate]{topsep=4pt,itemsep=2pt,parsep=0pt}
\hypersetup{colorlinks=true,linkcolor=blue!55!black,urlcolor=blue!55!black}

\titleformat{\section}{\large\bfseries}{\thesection.}{0.55em}{}[{\vspace{1pt}\titlerule[0.4pt]}]
\titleformat{\subsection}{\normalsize\bfseries}{\thesubsection.}{0.55em}{}
\titleformat{\subsubsection}{\small\bfseries\itshape}{\thesubsubsection.}{0.5em}{}

%--- ambienti teorema (solo enunciati, niente dimostrazioni)
\theoremstyle{plain}
\newtheorem{teorema}{Teorema}[section]
\newtheorem{proposizione}[teorema]{Proposizione}
\theoremstyle{definition}
\newtheorem{definizione}[teorema]{Definizione}
\newtheorem{esempio}[teorema]{Esempio}
\theoremstyle{remark}
\newtheorem*{osservazione}{Osservazione}
\newtheorem*{attenzione}{Nota numerica}

%--- macro acronimi (funzionano dentro e fuori math mode)
\newcommand{\UFR}{\ensuremath{\mathrm{UFR}}}
\newcommand{\LLP}{\ensuremath{\mathrm{LLP}}}
\newcommand{\CP}{\ensuremath{\mathrm{CP}}}
\newcommand{\CRA}{\ensuremath{\mathrm{CRA}}}
\newcommand{\Qmat}{\mathbf{Q}}
\newcommand{\Hmat}{\mathbf{H}}
\newcommand{\transp}{^{\!\top}}

%--- colore righe tabella
\definecolor{rowtint}{RGB}{243,247,253}

%--- box colorati
\tcbset{
  mybase/.style={
    colback=gray!5, colframe=gray!40,
    boxrule=0.4pt, arc=2pt,
    left=7pt, right=7pt, top=5pt, bottom=5pt,
    fonttitle=\small\bfseries
  }
}
\newtcolorbox{keybox}[1][Concetto chiave]{mybase,
  colframe=blue!40!black, colback=blue!4, title={#1}}
\newtcolorbox{algobox}[1][Algoritmo]{mybase,
  colframe=teal!50!black, colback=teal!4, title={#1}}
\newtcolorbox{warningbox}{mybase,
  colframe=orange!70!black, colback=orange!5}
\newtcolorbox{codebox}[1][Codice R]{mybase,
  colframe=green!40!black, colback=green!3, title={#1}}

% ============================================================
\title{%
  \textbf{Costruzione della Curva EIOPA Risk-Free}\\[6pt]
  \large Il metodo Smith--Wilson: un approccio variazionale\\[4pt]
  \normalsize La matrice di Wilson è definita positiva: un quadrato più l'integrale di un quadrato\\[2pt]
}
\author{Laboratorio di Calcolo Numerico --- Laurea Triennale in Matematica Applicata\\
  Universit\`a degli Studi di Verona\\[2pt]
  Anno Accademico 2026--2027}
\date{}

\begin{document}
\maketitle
\tableofcontents
\newpage

\noindent
Questa nota isola e sviluppa \textbf{un solo passaggio} della dispensa 02c (Sez.~4.2): la
dimostrazione che la matrice di Wilson $\Hmat$ è simmetrica definita positiva, ottenuta
riscrivendo la forma quadratica $\zeta\transp\Hmat\zeta$ come \textbf{un quadrato puntuale più
l'integrale di un quadrato}. Il contenuto matematico è identico a quello già presente in 02c
(righe 789--919) e riprodotto verbatim in 02e (Appendice A): qui però ogni passaggio è reso
esplicito --- gli integrali sono calcolati termine a termine, il passaggio dalla somma doppia
al quadrato è tracciato con un esempio a $m=2$ e uno a $m=3$, e ogni identità algebrica è
verificata anche numericamente. Le definizioni necessarie sono richiamate (non dimostrate) in
Sez.~\ref{sec:richiami}: per il contesto economico completo (perché si costruisce $\Hmat$, cosa
rappresenta $\zeta$) si rimanda a 02c, Sez.~2--4.1.

% ============================================================
\section{Richiami: nucleo di Wilson e forma quadratica}
\label{sec:richiami}
% ============================================================

Fissiamo qui, senza dimostrarle, le definizioni sulle quali poggia l'intera nota; sono le
stesse di 02c (Sez.~4.1--4.2), a cui si rimanda per la motivazione economica.

\begin{definizione}[Nucleo di Wilson, \cite{eiopa2025} sez.~9.7]
Per $\alpha>0$ fissato, il \textbf{nucleo di Wilson} è la funzione simmetrica
\[
  H(t,u) = \alpha\,\min(t,u) - e^{-\alpha\,\max(t,u)}\,\sinh\!\bigl(\alpha\,\min(t,u)\bigr),
  \qquad H(t,u)=H(u,t).
\]
\end{definizione}

\begin{definizione}[Curva di sconto Smith--Wilson]
Dati $m$ nodi di pagamento $0<u_1<\cdots<u_m$ e coefficienti $\zeta=(\zeta_1,\ldots,\zeta_m)\transp
\in\mathbb{R}^m$, la curva di sconto si scrive
\[
  P(t) = e^{-\omega t}\Bigl(1 + \sum_{j=1}^m \zeta_j\, H(t,u_j)\Bigr),
  \qquad \omega=\ln(1+\UFR).
\]
\end{definizione}

\begin{definizione}[Matrice di Wilson e forma quadratica]
La \textbf{matrice di Wilson} è $\Hmat = [H(u_i,u_j)]_{i,j=1}^m$; la sua \textbf{forma
quadratica associata} (l'\emph{energia} della curva) è
\[
  E(\zeta) = \zeta\transp\Hmat\,\zeta = \sum_{i=1}^m\sum_{j=1}^m \zeta_i\,\zeta_j\,H(u_i,u_j).
\]
\end{definizione}

\begin{keybox}[Perché conta]
Se $E(\zeta)=\zeta\transp\Hmat\zeta>0$ per ogni $\zeta\neq 0$, allora $\Hmat$ è simmetrica
definita positiva (SPD): il sistema lineare di calibrazione (02c, Sez.~4.3) ammette
un'\textbf{unica} soluzione, risolvibile con Cholesky, ed $E$ induce un vero prodotto scalare
$\langle x,y\rangle_\Hmat = x\transp\Hmat y$ su $\mathbb{R}^m$. Tutto il resto della dispensa 02c
--- convessità stretta del problema di minima energia, unicità della curva --- dipende da
questo unico fatto.
\end{keybox}

L'obiettivo di questa nota è dimostrare che $\Hmat$ è SPD usando \textbf{solo calcolo di
integrali}, tramite una rappresentazione ad hoc di $H$ come "prodotto scalare di funzioni".

% ============================================================
\section{Rappresentazione integrale del nucleo}
\label{sec:rappresentazione}
% ============================================================

\begin{definizione}[Le funzioni ausiliarie $\varphi_t$]
Per ogni $t>0$ si definisca
\[
  \varphi_t(x) =
  \begin{cases}
    1-e^{-\alpha(t-x)} & 0\le x\le t,\\[2pt]
    0 & x>t .
  \end{cases}
\]
Si osservi che $\varphi_t(0)=1-e^{-\alpha t}\in(0,1)$, che $\varphi_t$ è \textbf{continua} su
$[0,+\infty)$ (in $x=t$ entrambi i rami valgono $0$, senza salto), che $\varphi_t$ è
\textbf{crescente} su $[0,t]$ (derivata $\alpha e^{-\alpha(t-x)}>0$), e che il suo
\textbf{supporto} è esattamente $[0,t]$: fuori da quell'intervallo $\varphi_t\equiv 0$. Questa
proprietà di supporto compatto sarà l'ingrediente decisivo nella Sez.~\ref{sec:stretta}.
\end{definizione}

\begin{proposizione}[Rappresentazione integrale del nucleo di Wilson]
\label{prop:H-int}
Per ogni $s,t>0$ vale l'identità
\begin{equation}
  H(s,t) \;=\; \varphi_s(0)\,\varphi_t(0)
  \;+\; \alpha\!\int_0^{+\infty}\!\varphi_s(x)\,\varphi_t(x)\,dx .
  \label{eq:H-int}
\end{equation}
\end{proposizione}

Dimostriamo la~\eqref{eq:H-int} con tutti i passaggi intermedi. Per la simmetria di ambo i
membri in $s,t$, non è restrittivo assumere $s\le t$.

\paragraph{Passo 0: dove vive l'integrale.} Il prodotto $\varphi_s(x)\varphi_t(x)$ è nullo non
appena $x>s$ (perché $\varphi_s(x)=0$ lì, essendo $s\le t$), quindi
\[
  \int_0^{+\infty}\varphi_s(x)\varphi_t(x)\,dx = \int_0^{s}\varphi_s(x)\varphi_t(x)\,dx ,
\]
un integrale su un intervallo \textbf{limitato} di una funzione \textbf{continua} (prodotto di
continue): nessun problema di convergenza.

\paragraph{Passo 1: il prodotto $\varphi_s(x)\varphi_t(x)$ sull'intervallo $[0,s]$.} Su questo
intervallo $x\le s\le t$, quindi \emph{entrambi} i rami "interni" delle definizioni si
applicano: $\varphi_s(x)=1-e^{-\alpha(s-x)}$ e $\varphi_t(x)=1-e^{-\alpha(t-x)}$. Svolgendo il
prodotto di questi due binomi,
\[
  \varphi_s(x)\varphi_t(x)
  = \bigl(1-e^{-\alpha(s-x)}\bigr)\bigl(1-e^{-\alpha(t-x)}\bigr)
  = 1 - e^{-\alpha(s-x)} - e^{-\alpha(t-x)} + e^{-\alpha(s-x)}e^{-\alpha(t-x)} ,
\]
e osservando che $e^{-\alpha(s-x)}e^{-\alpha(t-x)}=e^{-\alpha(s+t-2x)}$,
\begin{equation}
  \varphi_s(x)\,\varphi_t(x) = 1 - e^{-\alpha(s-x)} - e^{-\alpha(t-x)} + e^{-\alpha(s+t-2x)},
  \qquad x\in[0,s].
  \label{eq:prodotto-espanso}
\end{equation}
Quattro addendi: una costante e tre esponenziali in $x$. Li integriamo separatamente.

\paragraph{Passo 2: i quattro integrali, uno alla volta.} Su $[0,s]$:
\begin{align}
  \int_0^s 1\,dx &= s, \label{eq:int1}\\[4pt]
  \int_0^s e^{-\alpha(s-x)}\,dx
  &= \Bigl[\tfrac{1}{\alpha}e^{-\alpha(s-x)}\Bigr]_0^s
   = \tfrac{1}{\alpha}\bigl(1-e^{-\alpha s}\bigr), \label{eq:int2}\\[4pt]
  \int_0^s e^{-\alpha(t-x)}\,dx
  &= \Bigl[\tfrac{1}{\alpha}e^{-\alpha(t-x)}\Bigr]_0^s
   = \tfrac{1}{\alpha}\bigl(e^{-\alpha(t-s)}-e^{-\alpha t}\bigr), \label{eq:int3}\\[4pt]
  \int_0^s e^{-\alpha(s+t-2x)}\,dx
  &= \Bigl[\tfrac{1}{2\alpha}e^{-\alpha(s+t-2x)}\Bigr]_0^s
   = \tfrac{1}{2\alpha}\bigl(e^{-\alpha(t-s)}-e^{-\alpha(s+t)}\bigr). \label{eq:int4}
\end{align}
In~\eqref{eq:int2}--\eqref{eq:int4} si è usato $\frac{d}{dx}e^{\pm\beta x}=\pm\beta\,e^{\pm\beta x}$
con $\beta=\alpha>0$ (in~\eqref{eq:int4}, $\beta=2\alpha$): ogni primitiva è verificabile per
derivazione diretta. Si noti in~\eqref{eq:int3} che l'esponente $-\alpha(t-x)$, valutato in
$x=s$, dà $-\alpha(t-s)$ e \emph{non} $-\alpha(s-t)$: $t-s\ge 0$ perché $s\le t$, quindi
$e^{-\alpha(t-s)}\in(0,1]$ come deve essere.

\paragraph{Passo 3: assemblare l'integrale, moltiplicare per $\alpha$.} Sommando
\eqref{eq:int1}--\eqref{eq:int4} con i segni di~\eqref{eq:prodotto-espanso}
($+,-,-,+$) e moltiplicando per $\alpha$ (i fattori $\tfrac1\alpha$ e $\tfrac1{2\alpha}$ si
semplificano):
\begin{equation}
  \alpha\!\int_0^{s}\!\varphi_s(x)\varphi_t(x)\,dx
  = \underbrace{\alpha s}_{\text{da }\eqref{eq:int1}}
    \;\underbrace{-\bigl(1-e^{-\alpha s}\bigr)}_{\text{da }\eqref{eq:int2}}
    \;\underbrace{-\bigl(e^{-\alpha(t-s)}-e^{-\alpha t}\bigr)}_{\text{da }\eqref{eq:int3}}
    \;\underbrace{+\tfrac12\bigl(e^{-\alpha(t-s)}-e^{-\alpha(s+t)}\bigr)}_{\text{da }\eqref{eq:int4}} .
  \label{eq:somma-integrale}
\end{equation}

\paragraph{Passo 4: il termine puntuale.} Il primo addendo di~\eqref{eq:H-int} si calcola in
un colpo solo (nessun integrale, solo prodotto di due valori):
\begin{equation}
  \varphi_s(0)\,\varphi_t(0)
  = \bigl(1-e^{-\alpha s}\bigr)\bigl(1-e^{-\alpha t}\bigr)
  = 1 - e^{-\alpha s} - e^{-\alpha t} + e^{-\alpha(s+t)} .
  \label{eq:punto}
\end{equation}

\paragraph{Passo 5: cancellazione termine a termine.} Sommiamo~\eqref{eq:somma-integrale} e
\eqref{eq:punto}. Per non lasciare nulla all'"si verifica che", raccogliamo \emph{ogni} tipo
di termine in una riga a sé:

\begin{center}
\renewcommand{\arraystretch}{1.3}
\begin{tabular}{lcccccc}
\toprule
\textbf{Origine} & $\alpha s$ & $1$ & $e^{-\alpha s}$ & $e^{-\alpha t}$
  & $e^{-\alpha(t-s)}$ & $e^{-\alpha(s+t)}$ \\
\midrule
da $\eqref{eq:int1}$, coeff. $+$ & $+1$ & & & & & \\
da $\eqref{eq:int2}$, coeff. $-$ & & $-1$ & $+1$ & & & \\
da $\eqref{eq:int3}$, coeff. $-$ & & & & $+1$ & $-1$ & \\
da $\eqref{eq:int4}$, coeff. $+\tfrac12$ & & & & & $+\tfrac12$ & $-\tfrac12$ \\
da $\eqref{eq:punto}$ & & $+1$ & $-1$ & $-1$ & & $+1$ \\
\midrule
\textbf{somma delle colonne} & $\mathbf{+1}$ & $0$ & $0$ & $0$
  & $\mathbf{-\tfrac12}$ & $\mathbf{+\tfrac12}$ \\
\bottomrule
\end{tabular}
\end{center}

Le colonne "$1$", "$e^{-\alpha s}$" ed "$e^{-\alpha t}$" si annullano \textbf{esattamente}:
sopravvivono solo $\alpha s$, $-\tfrac12 e^{-\alpha(t-s)}$ e $+\tfrac12 e^{-\alpha(s+t)}$. Quindi
\begin{equation}
  H(s,t) \;\overset{?}{=}\; \varphi_s(0)\varphi_t(0)+\alpha\!\int_0^{+\infty}\!\varphi_s\varphi_t\,dx
  \;=\; \alpha s - \tfrac12 e^{-\alpha(t-s)} + \tfrac12 e^{-\alpha(s+t)} .
  \label{eq:dopo-cancellazione}
\end{equation}

\paragraph{Passo 6: riconoscere il seno iperbolico.} Resta da mostrare che il membro destro di
\eqref{eq:dopo-cancellazione} coincide con la definizione di $H(s,t)$. Raccogliamo
$e^{-\alpha t}$ dagli ultimi due termini:
\[
  -\tfrac12 e^{-\alpha(t-s)} + \tfrac12 e^{-\alpha(s+t)}
  = -\tfrac12 e^{-\alpha t}e^{\alpha s} + \tfrac12 e^{-\alpha t}e^{-\alpha s}
  = -e^{-\alpha t}\cdot\frac{e^{\alpha s}-e^{-\alpha s}}{2}
  = -e^{-\alpha t}\,\sinh(\alpha s),
\]
usando la definizione $\sinh(z)=\tfrac12(e^{z}-e^{-z})$. Sostituendo in~\eqref{eq:dopo-cancellazione},
\begin{align*}
  \varphi_s(0)\varphi_t(0)+\alpha\!\int_0^{+\infty}\!\varphi_s\varphi_t\,dx
  &= \alpha s - e^{-\alpha t}\sinh(\alpha s)\\
  &= \alpha\min(s,t) - e^{-\alpha\max(s,t)}\sinh\!\bigl(\alpha\min(s,t)\bigr) \;=\; H(s,t),
\end{align*}
dove l'ultimo passaggio usa $s\le t$, cioè $s=\min(s,t)$ e $t=\max(s,t)$. Questo dimostra
\eqref{eq:H-int}. $\blacksquare$

\begin{osservazione}[Per approfondire]
La rappresentazione~\eqref{eq:H-int} è un caso particolare di una tecnica generale per
costruire nuclei definiti positivi come funzioni di Green di un operatore differenziale: si
veda Wahba~\cite{wahba1990}, cap.~1, per la costruzione analoga (con un operatore del
second'ordine) alla base delle spline cubiche di lisciamento.
\end{osservazione}

\begin{attenzione}
Verifichiamo~\eqref{eq:H-int} con numeri, per $\alpha=0.0736$ (valore ufficiale calibrato su
dicembre 2025, 02c Sez.~5.3) e $s=1,\,t=5$. Termine puntuale:
$\varphi_1(0)\varphi_5(0) = (1-e^{-0.0736})(1-e^{-0.368}) = 0.070957\times 0.307883 = 0.021846$.
Termine integrale (calcolato numericamente su $[0,1]$, dove il prodotto è supportato):
$\alpha\int_0^{1}\varphi_1(x)\varphi_5(x)\,dx = 0.000768$. Somma: $0.021846+0.000768=0.022614$.
Dalla formula chiusa, $H(1,5)=\alpha\cdot 1 - e^{-\alpha\cdot 5}\sinh(\alpha\cdot 1)
= 0.0736 - 0.692117\times 0.073666 = 0.022614$: i due membri di~\eqref{eq:H-int} coincidono
(la differenza è dell'ordine di $10^{-17}$, puro arrotondamento in doppia precisione). Il
punto pedagogico è che l'identità va verificata con un calcolo, non presa per fede.
\end{attenzione}

% ============================================================
\section{Semipositività: un quadrato più l'integrale di un quadrato}
\label{sec:semipositivita}
% ============================================================

\begin{proposizione}[La matrice di Wilson è semidefinita positiva]
\label{prop:H-spd}
Per ogni scelta di nodi $0<u_1<\cdots<u_m$ e per ogni $\zeta\in\mathbb{R}^m$,
$\zeta\transp\Hmat\zeta\ge 0$.
\end{proposizione}

\paragraph{Passo 1: sostituire la rappresentazione integrale in ogni entrata.} Per definizione,
\[
  \zeta\transp\Hmat\zeta = \sum_{i=1}^m\sum_{j=1}^m \zeta_i\,\zeta_j\,H(u_i,u_j).
\]
Sostituendo~\eqref{eq:H-int} in \emph{ciascuno} degli $m^2$ addendi $H(u_i,u_j)$,
\[
  \zeta\transp\Hmat\zeta
  = \sum_{i=1}^m\sum_{j=1}^m \zeta_i\zeta_j
    \Bigl[\varphi_{u_i}(0)\varphi_{u_j}(0) + \alpha\!\int_0^{+\infty}\!\varphi_{u_i}(x)\varphi_{u_j}(x)\,dx\Bigr] .
\]

\paragraph{Passo 2: distribuire la somma doppia sulle due parti.} La somma doppia ha
\textbf{esattamente} $m^2$ addendi (un numero finito): possiamo quindi spezzarla nella somma
delle due parti (puntuale e integrale) senza nessuna ipotesi di convergenza o di scambio
somma/integrale in senso analitico --- non serve altro che l'additività di una somma finita:
\[
  \zeta\transp\Hmat\zeta
  = \underbrace{\sum_{i,j}\zeta_i\zeta_j\,\varphi_{u_i}(0)\varphi_{u_j}(0)}_{=:A}
  \;+\; \alpha\underbrace{\sum_{i,j}\zeta_i\zeta_j\!\int_0^{+\infty}\!\varphi_{u_i}(x)\varphi_{u_j}(x)\,dx}_{=:B(x)\text{, poi integrato}} .
\]
Nel secondo blocco, per lo stesso motivo (somma finita di $m^2$ integrali), possiamo portare
$\sum_{i,j}$ \textbf{dentro} l'integrale:
\[
  \alpha\sum_{i,j}\zeta_i\zeta_j\!\int_0^{+\infty}\!\varphi_{u_i}(x)\varphi_{u_j}(x)\,dx
  = \alpha\!\int_0^{+\infty}\!\Bigl(\sum_{i,j}\zeta_i\zeta_j\,\varphi_{u_i}(x)\varphi_{u_j}(x)\Bigr)dx .
\]

\paragraph{Passo 3: il quadrato di una somma, prima a $m=2$ poi in generale.} Poniamo
$h(x):=\sum_{k=1}^m \zeta_k\,\varphi_{u_k}(x)$ (per $x=0$: $h(0)=\sum_k\zeta_k\varphi_{u_k}(0)$).
L'identità chiave, per qualunque scelta di $m$ numeri reali $a_1,\ldots,a_m$, è
\[
  \Bigl(\sum_{k=1}^m a_k\Bigr)^{\!2} = \sum_{i=1}^m\sum_{j=1}^m a_i a_j .
\]
Verifichiamola prima nel caso più piccolo non banale, $m=2$: sviluppando il quadrato di un
binomio, $(a_1+a_2)^2 = a_1^2+2a_1a_2+a_2^2$; d'altra parte la somma doppia con $i,j\in\{1,2\}$
ha \textbf{quattro} addendi (non due),
\[
  \sum_{i=1}^2\sum_{j=1}^2 a_ia_j = a_1a_1+a_1a_2+a_2a_1+a_2a_2 = a_1^2+2a_1a_2+a_2^2,
\]
perché $a_1a_2$ e $a_2a_1$ sono lo \textbf{stesso numero} e vengono entrambi contati (la somma
doppia scorre su \emph{coppie ordinate}): la loro somma ricostruisce il doppio prodotto
$2a_1a_2$. I due membri coincidono. Il caso generale è identico, solo con più addendi: ogni
coppia non ordinata $\{i,j\}$ con $i\neq j$ compare due volte in $\sum_{i,j}a_ia_j$ (come
$(i,j)$ e come $(j,i)$), producendo $2a_ia_j$, mentre ogni $i=j$ compare una volta sola,
producendo $a_i^2$ --- esattamente lo sviluppo del quadrato della somma.

Applicando questa identità con $a_k=\zeta_k\varphi_{u_k}(0)$ al blocco $A$ del Passo 2, e con
$a_k=\zeta_k\varphi_{u_k}(x)$ (per $x$ fissato) al blocco sotto integrale:
\[
  A = \sum_{i,j}\zeta_i\zeta_j\varphi_{u_i}(0)\varphi_{u_j}(0)
    = \Bigl(\sum_k \zeta_k\varphi_{u_k}(0)\Bigr)^{\!2} = h(0)^2,
  \qquad
  \sum_{i,j}\zeta_i\zeta_j\varphi_{u_i}(x)\varphi_{u_j}(x) = h(x)^2 .
\]

\paragraph{Passo 4: conclusione.} Sostituendo nel Passo 2,
\begin{equation}
  \zeta\transp\Hmat\,\zeta = h(0)^2 + \alpha\!\int_0^{+\infty}\! h(x)^2\,dx .
  \label{eq:energia-int}
\end{equation}
Il membro destro è la somma di \textbf{due quantità non negative} ($h(0)^2\ge 0$ perché è un
quadrato; $\alpha\int h^2\ge 0$ perché $\alpha>0$ e l'integrale di una funzione non negativa è
non negativo): dunque $\zeta\transp\Hmat\zeta\ge 0$ per ogni $\zeta$, cioè $\Hmat$ è semidefinita
positiva. $\blacksquare$

\begin{osservazione}[I termini misti $i\neq j$ non spariscono]
Un errore frequente è pensare che il passaggio "somma doppia $\to$ quadrato" scarti i prodotti
incrociati $\zeta_i\zeta_j$ con $i\neq j$. È il contrario: come mostrato nel Passo~3,
$\bigl(\sum_k a_k\bigr)^2$ \emph{contiene già tutti} i termini $a_ia_j$, compresi quelli misti
(con molteplicità 2 ciascuno). La~\eqref{eq:energia-int} è un'\textbf{identità esatta}, non
un'approssimazione che ignora le correlazioni fra nodi diversi.
\end{osservazione}

\begin{attenzione}[Verifica numerica completa di~\eqref{eq:energia-int}]
Prendiamo $m=3$ nodi $u=(1,5,10)$, $\alpha=0.0736$ (valore ufficiale, 02c Sez.~5.3) e un vettore
\textbf{arbitrario} $\zeta=(2,-1,0.5)\transp$ (nessun significato finanziario: serve solo a
verificare l'identità con numeri concreti, non è legato all'esempio della Sez.~\ref{sec:stretta}).
Dalla formula chiusa di $H$:
\[
  \Hmat = \begin{pmatrix}
    0.005161 & 0.022614 & 0.038312\\
    0.022614 & 0.107513 & 0.187713\\
    0.038312 & 0.187713 & 0.350733
  \end{pmatrix},
  \qquad
  \zeta\transp\Hmat\zeta = 0.0142935 .
\]
Dall'altro lato, $h(0)=\sum_k\zeta_k\varphi_{u_k}(0) = 2(0.070957)+(-1)(0.307883)+0.5(0.520974)
= 0.094518$, quindi $h(0)^2=0.0089336$. Integrando numericamente $h(x)^2$ su $[0,10]$ (oltre
$x=10$, $h\equiv 0$ perché tutti e tre i $\varphi_{u_k}$ sono fuori supporto) con una griglia
fine si ottiene $\int_0^{10}h(x)^2\,dx = 0.0728245$, dunque $\alpha\int h^2 = 0.0053599$. Somma:
$0.0089336+0.0053599=0.0142935$, che coincide con $\zeta\transp\Hmat\zeta$ calcolato sopra fino
alla settima cifra significativa (la differenza residua, $\sim\!10^{-15}$, è arrotondamento
numerico dell'integrazione, non un errore dell'identità).
\end{attenzione}

% ============================================================
\section{Positività stretta: induzione discendente}
\label{sec:stretta}
% ============================================================

La Proposizione~\ref{prop:H-spd} dà $\zeta\transp\Hmat\zeta\ge 0$: per concludere che $\Hmat$ è
\textbf{definita} (non solo semidefinita) positiva serve mostrare che l'uguaglianza a zero
forza $\zeta=0$.

\begin{proposizione}[La matrice di Wilson è SPD]
\label{prop:H-spd-stretta}
Se i nodi $0<u_1<\cdots<u_m$ sono \textbf{distinti}, allora $\zeta\transp\Hmat\zeta=0$ implica
$\zeta=0$. Insieme alla Proposizione~\ref{prop:H-spd}, $\Hmat$ è simmetrica definita positiva.
\end{proposizione}

\paragraph{Passo 1: dalla somma nulla ai due addendi nulli separatamente.} Supponiamo
$\zeta\transp\Hmat\zeta=0$. Per~\eqref{eq:energia-int} questo significa
$h(0)^2+\alpha\int_0^{+\infty}h(x)^2\,dx=0$, somma di \textbf{due} termini non negativi. Una
somma di numeri non negativi è zero se e solo se \emph{ciascun} termine è zero (se uno dei due
fosse strettamente positivo, la somma lo sarebbe anch'essa): dunque
\[
  h(0)^2=0 \quad\Longrightarrow\quad h(0)=0,
  \qquad\qquad
  \int_0^{+\infty}h(x)^2\,dx = 0 .
\]

\paragraph{Passo 2: dall'integrale nullo a $h\equiv 0$.} Ogni $\varphi_{u_k}$ è continua (Sez.~
\ref{sec:rappresentazione}), quindi $h=\sum_k\zeta_k\varphi_{u_k}$, somma finita di funzioni
continue, è continua. Una funzione continua e non negativa ($h^2\ge 0$) il cui integrale su
$[0,+\infty)$ è nullo deve essere identicamente nulla: se fosse $h(x_0)^2=\varepsilon>0$ in un
punto $x_0$, per continuità $h^2>\varepsilon/2$ su un intorno di $x_0$ di lunghezza positiva
$\delta$, e allora $\int h^2 \ge (\varepsilon/2)\delta>0$, contraddizione. Dunque
\begin{equation}
  h(x) = \sum_{k=1}^m \zeta_k\,\varphi_{u_k}(x) = 0 \qquad \text{per \emph{ogni} } x\ge 0.
  \label{eq:h-zero}
\end{equation}

\paragraph{Passo 3: da $h\equiv 0$ a $\zeta=0$, isolando un coefficiente alla volta.} Qui entra
in gioco l'unica proprietà di $\varphi_{u_k}$ non ancora usata: il \textbf{supporto compatto}
$[0,u_k]$. Poniamo $u_0:=0$ e procediamo per \textbf{induzione discendente} su
$k=m,m-1,\ldots,1$, mostrando $\zeta_k=0$ un nodo alla volta, dal più lontano al più vicino.

\emph{Caso base $k=m$.} Fissiamo $x$ nell'intervallo \textbf{aperto} $(u_{m-1},u_m)$. Per ogni
$j<m$ si ha $u_j\le u_{m-1}<x$, cioè $x$ è fuori dal supporto $[0,u_j]$ di $\varphi_{u_j}$:
$\varphi_{u_j}(x)=0$. Nella somma~\eqref{eq:h-zero} sopravvive quindi solo il termine $j=m$:
\[
  0 = h(x) = \zeta_m\,\varphi_{u_m}(x) \qquad \forall\, x\in(u_{m-1},u_m).
\]
La funzione $\varphi_{u_m}(x)=1-e^{-\alpha(u_m-x)}$ si annulla \emph{solo} in $x=u_m$, punto
\textbf{escluso} dall'intervallo aperto; ad esempio nel punto medio
$\bar x=(u_{m-1}+u_m)/2$ vale $\varphi_{u_m}(\bar x)=1-e^{-\alpha(u_m-\bar x)}\neq 0$ (è
strettamente positiva, essendo $\varphi_{u_m}$ crescente da $0$ a $\varphi_{u_m}(u_m^-)$ sul suo
supporto). Da $\zeta_m\,\varphi_{u_m}(\bar x)=0$ con $\varphi_{u_m}(\bar x)\neq 0$ segue
$\zeta_m=0$.

\emph{Ipotesi induttiva.} Sia $1\le k<m$ e supponiamo già dimostrato $\zeta_{k+1}=\cdots=\zeta_m=0$.

\emph{Passo induttivo.} Fissiamo $x\in(u_{k-1},u_k)$ (con $u_0=0$ se $k=1$). Per $j<k$, come nel
caso base, $u_j\le u_{k-1}<x$ dà $\varphi_{u_j}(x)=0$. Per $j>k$, invece, il coefficiente
$\zeta_j$ è \textbf{già nullo per ipotesi induttiva}: il termine $\zeta_j\varphi_{u_j}(x)$ si
annulla per via del coefficiente, \emph{indipendentemente} dal fatto che $\varphi_{u_j}(x)$
possa essere non nulla in quel punto (il nodo $u_j$ è ancora "avanti" rispetto a $x$, quindi
$x$ è dentro il suo supporto). Resta un solo termine:
\[
  0 = h(x) = \zeta_k\,\varphi_{u_k}(x) \qquad \forall\, x\in(u_{k-1},u_k),
\]
e lo stesso argomento del caso base (applicato ora al nodo $u_k$: $\varphi_{u_k}$ non è
identicamente nulla sull'intervallo aperto $(u_{k-1},u_k)$) dà $\zeta_k=0$.

Per induzione, $\zeta_k=0$ per ogni $k=m,m-1,\ldots,1$, cioè $\zeta=0$. $\blacksquare$

\begin{osservazione}[Perché servono nodi distinti]
L'ipotesi $u_1<\cdots<u_m$ (nodi \textbf{distinti}, non solo $\le$) è ciò che rende ogni
intervallo $(u_{k-1},u_k)$ non vuoto: se due nodi coincidessero, l'intervallo corrispondente
degenererebbe a un punto (o sparirebbe) e il Passo~3 non avrebbe più un punto interno $\bar x$
su cui valutare $h$ e isolare $\zeta_k$.
\end{osservazione}

\begin{osservazione}[Per approfondire: una dimostrazione probabilistica alternativa]
Una via del tutto diversa per la stessa Proposizione~\ref{prop:H-spd-stretta}, per via
probabilistica (total positivity del nucleo di un processo di Ornstein--Uhlenbeck e formula di
Cauchy--Binet), si trova in Lager\aa s e Lindholm~\cite{lageras2016}, sez.~3 e~5; qui non la
sviluppiamo, perché l'obiettivo di questa nota è restare autosufficienti con i soli integrali
già introdotti in Sez.~\ref{sec:rappresentazione}.
\end{osservazione}

\paragraph{Esempio: l'induzione con $m=3$ nodi.} Rendiamo concreto il Passo~3 con
$u=(u_1,u_2,u_3)=(1,5,10)$ e $\alpha=0.0736$. \textbf{Attenzione}: qui $\zeta$ non è più il
vettore arbitrario $(2,-1,0.5)$ della Sez.~\ref{sec:semipositivita} (che serviva solo a
verificare l'identità~\eqref{eq:energia-int}); l'ipotesi di questa sezione è invece
$h\equiv 0$, cioè esattamente la situazione (puramente ipotetica, usata per dimostrare
l'implicazione) in cui si vuole mostrare che l'\emph{unico} $\zeta$ compatibile è $\zeta=0$.
La tabella elenca, intervallo per intervallo, quali funzioni ausiliarie sono dentro il proprio
supporto e quale coefficiente l'argomento isola:

\begin{center}
\small
\renewcommand{\arraystretch}{1.35}
\setlength{\tabcolsep}{4.5pt}
\begin{tabular}{lccclc}
\toprule
\textbf{Intervallo} & $\varphi_{u_1}$ & $\varphi_{u_2}$ & $\varphi_{u_3}$
  & \textbf{$h(x)$ si riduce a} & \textbf{concl.}\\
\midrule
$(5,10)$ & fuori supp. & fuori supp. & in supp.
  & $\zeta_3\,\varphi_{u_3}(x)$ & $\zeta_3=0$\\
$(1,5)$ & fuori supp. & in supp. & $\zeta_3=0$ già noto
  & $\zeta_2\,\varphi_{u_2}(x)$ & $\zeta_2=0$\\
$(0,1)$ & in supp. & $\zeta_2=0$ già noto & $\zeta_3=0$ già noto
  & $\zeta_1\,\varphi_{u_1}(x)$ & $\zeta_1=0$\\
\bottomrule
\end{tabular}

\smallskip
{\footnotesize supp.\ $=$ supporto: $\varphi_{u_1}$ su $[0,1]$, $\varphi_{u_2}$ su $[0,5]$,
$\varphi_{u_3}$ su $[0,10]$.}
\end{center}

Come testimone concreto che $\varphi_{u_3}$ non è identicamente nulla su $(5,10)$: nel punto
medio $\bar x=7.5$, $\varphi_{u_3}(7.5)=1-e^{-0.0736\times 2.5}=1-e^{-0.184}=0.1681\neq 0$ ---
il valore usato implicitamente nel caso base della dimostrazione generale. L'ordine
"dal nodo più lontano al più vicino" non è arbitrario: è ciò che garantisce che, quando si
esamina l'intervallo $(u_{k-1},u_k)$, tutti i coefficienti $\zeta_j$ con $j>k$ siano \emph{già}
stati azzerati dal passo precedente.

% ============================================================
\section{Confronto e chiusura}
\label{sec:confronto}
% ============================================================

Questa nota espande, senza modificarlo nella sostanza, il paragrafo "Semipositività: la forma
quadratica è un quadrato più l'integrale di un quadrato" già presente in 02c (Sez.~4.2, righe
841--863) e riprodotto verbatim in 02e (Appendice A, righe 1668--1692). La dimostrazione resta
\textbf{la stessa} nelle sue tre proposizioni cardine (rappresentazione integrale, quadrato più
integrale di un quadrato, induzione discendente): qui ogni passaggio algebrico che in 02c/02e è
compresso in una o due righe è stato scomposto in passi numerati e ogni identità è stata
verificata anche con numeri concreti ($\alpha=0.0736$, i nodi $u=(1,5,10)$).

\begin{osservazione}[Relazione con 02c e 02e]
02c dimostra la positiva definitezza di $\Hmat$ con \emph{esattamente} questo argomento (funzioni
ausiliarie $\varphi_t$ più induzione), in forma compatta, all'interno della trattazione
economica completa del metodo Smith--Wilson. 02e ottiene lo stesso risultato per una via
alternativa --- il funzionale di energia che Smith e Wilson minimizzano davvero induce un
prodotto scalare di cui $W(t,u)$ è il rappresentante della valutazione puntuale, rendendo
$[W(u_i,u_j)]$ letteralmente una matrice di Gram --- e conserva la dimostrazione elementare qui
approfondita nella propria Appendice A, per chi preferisce un argomento che non richiede
l'apparato del funzionale di energia. Questa nota non introduce una terza dimostrazione: è una
lente di ingrandimento sulla prima.
\end{osservazione}

% ============================================================
\renewcommand{\refname}{Riferimenti bibliografici}
\begin{thebibliography}{9}

\bibitem{eiopa2025}
  \textbf{EIOPA} (2025).
  \textit{Technical documentation of the methodology to derive EIOPA's risk-free interest rate
  term structures}, EIOPA-BoS-25-599.
  Sezione~9 (Extrapolation and interpolation) per il metodo Smith--Wilson.

\bibitem{fasshauer2007}
  \textbf{Fasshauer, G.E.} (2007).
  \textit{Meshfree Approximation Methods with MATLAB}. World Scientific.
  Interpolazione con kernel definiti positivi: l'interpolante è combinazione lineare delle
  sezioni del kernel e i coefficienti risolvono il sistema con la matrice di Gram SPD.

\bibitem{lageras2016}
  \textbf{Lager\aa s, A., Lindholm, M.} (2016).
  Issues with the Smith--Wilson method.
  \textit{Insurance: Mathematics and Economics}, 68, 32--41.
  Dimostrazione alternativa della positiva definitezza del nucleo via total positivity e
  formula di Cauchy--Binet (sez.~3 e~5).

\bibitem{wahba1990}
  \textbf{Wahba, G.} (1990).
  \textit{Spline Models for Observational Data}.
  CBMS-NSF Regional Conference Series in Applied Mathematics, Vol.~59. SIAM.
  Cap.~1: costruzione di nuclei definiti positivi come funzioni di Green di un operatore
  differenziale --- la tecnica generale di cui la Proposizione~\ref{prop:H-int} è un caso
  particolare.

\end{thebibliography}

\end{document}
